\documentclass[article,12pt]{article}
\usepackage{amsmath,amssymb,amsthm}
\usepackage{graphicx}
\usepackage{hyperref}
\usepackage{listings}
\usepackage{url}
\usepackage{cite}
\usepackage{geometry}
\geometry{margin=1in}

\title{LLVM Transactional Memory Plugin:\\Automatic Instrumentation for Software Transactional Memory}
\author{}
\date{\today}

\begin{document}

\maketitle

\begin{abstract}
We present an LLVM compiler pass that automatically instruments C++ code for Software Transactional Memory (STM). Using source-code annotations, developers mark global variables and functions as transactional, and our pass automatically inserts begin/commit infrastructure, read/write logging, and conflict detection. The plugin supports multiple STM backends including TinySTM, SingleGlobalLock, TL2, and SwissTM. We evaluate the plugin with a suite of benchmarks including banking transactions, data structures, STAMP, and YCSB workloads.
\end{abstract}

\section{Introduction}

Software Transactional Memory (STM) provides a promising abstraction for concurrent programming, allowing developers to specify atomic blocks of code that either complete atomically or abort and retry. However, manual STM implementation is error-prone and breaks abstraction.

We present an LLVM compiler pass that automates STM instrumentation:

\begin{itemize}
\item \textbf{Global variables} marked with \texttt{\_\_attribute\_\_((annotate("tm")))} are managed by the STM runtime
\item \textbf{Functions} marked with \texttt{\_\_attribute\_\_((annotate("transaction")))} are executed as transactions
\item The pass automatically inserts:
  \begin{itemize}
  \item Transaction begin/commit hooks
  \item Nested transaction detection via counter
  \item sigsetjmp/siglongjmp for automatic retry
  \item Read/write instrumentation for memory accesses
  \end{itemize}
\end{itemize}

%\newpage

\section{Implementation}\label{sec:implementation}

Our LLVM plugin consists of two compiler passes that operate in the standard LLVM Pass Manager framework:

\subsection{TMGlobalInitPass (Module-Level)}

This pass performs one-time initialization for the entire compilation unit:

\begin{enumerate}
\item \textbf{Declare Runtime Functions}: Insert declarations for:
  \begin{itemize}
  \item \texttt{tm\_init()}: Global initialization
  \item \texttt{tm\_exit()}: Global cleanup
  \item \texttt{tm\_init\_thread()}: Per-thread initialization
  \item \texttt{tm\_exit\_thread()}: Per-thread cleanup
  \end{itemize}
\item \textbf{Create Symbol Tables}: For TM-managed globals, create:
  \begin{itemize}
  \item \texttt{tm\_symbol\_names[]}: Array of global names
  \item \texttt{tm\_symbol\_addresses[]}: Array of global addresses
  \item \texttt{tm\_symbol\_count}: Number of TM globals
  \end{itemize}
\item \textbf{Insert Main Hooks}: In \texttt{main()}:
  \begin{itemize}
  \item Call \texttt{tm\_init()} at start
  \item Call \texttt{tm\_init\_thread()} at start
  \item Call \texttt{tm\_exit\_thread()} before each return
  \item Call \texttt{tm\_exit()} before each return
  \end{itemize}
\item \textbf{Thread Entry Points}: For functions using TM globals:
  \begin{itemize}
  \item Insert \texttt{tm\_init\_thread()} at entry
  \item Insert \texttt{tm\_exit\_thread()} before returns
  \end{itemize}
\end{enumerate}

\subsection{TMInstrumentPass (Function-Level)}

This pass instruments each transaction function:

\begin{enumerate}
\item \textbf{Detect Transaction Functions}: Check for \texttt{\_\_attribute\_\_((annotate("transaction")))}
\item \textbf{Nested Transaction Handling}:
  \begin{itemize}
  \item Create thread-local \texttt{tm\_nested\_call\_counter}
  \item Create thread-local \texttt{tm\_jmpbuf} (256 bytes)
  \item Create thread-local \texttt{tm\_longjmp\_ret}
  \end{itemize}
\item \textbf{Entry Instrumentation}:
  \begin{itemize}
  \item Load counter, check if 0 (outermost transaction)
  \item If outermost (counter=0): Call sigsetjmp, store result, call \texttt{tm\_begin()}
  \item If nested (counter>0 or retry): Skip to nested path
  \end{itemize}
\item \textbf{Exit Instrumentation}:
  \begin{itemize}
  \item Decrement counter
  \item If counter reaches 0: Call \texttt{tm\_end()}
  \end{itemize}
\item \textbf{Memory Instrumentation}:
  \begin{itemize}
  \item All loads $\rightarrow$ call \texttt{tm\_read\_TYPE(ptr)}
  \item All stores $\rightarrow$ call \texttt{tm\_write\_TYPE(ptr, value)}
  \item Supported types: i8, i16, i32, i64, float, double, pointers
  \end{itemize}
\end{enumerate}

Figure~\ref{fig:flow} shows the instrumentation flow:

\begin{figure}[h]
\centering
\fbox{\parbox{0.95\linewidth}{
\begin{center}
\texttt{// Original}\\[2mm]
\texttt{TX void transfer(Account* from, Account* to, int amount) \{}\\
\texttt{\  from->balance -= amount;}\\
\texttt{\  to->balance += amount;}\\
\texttt{\}}\\[4mm]
\hrulefill\\[4mm]
\texttt{// Instrumented}\\[2mm]
\texttt{// Entry}\\
\texttt{if (counter == 0) \{}\\
\texttt{\  sigsetjmp(jmpbuf, 0);}\\
\texttt{\  tm\_begin();}\\
\texttt{\}}\\[2mm]
\texttt{// Memory ops}\\
\texttt{tm\_write\_i4(\&from->balance, tm\_read\_i4(\&from->balance) - amount);}\\
\texttt{tm\_write\_i4(\&to->balance, tm\_read\_i4(\&to->balance) + amount);}\\[2mm]
\texttt{// Exit}\\
\texttt{tm\_end();}
\end{center}
}}
\caption{Transaction instrumentation flow}
\label{fig:flow}
\end{figure}

\subsection{Nested Transaction Detection}

For correctly handling nested transactions, we use a thread-local counter:

\begin{verbatim}
entry:
  %counter = load tm_nested_call_counter
  %is_outer = icmp eq %counter, 0
  br %is_outer, label %outer, label %nested

outer:
  %setjmp_ret = call sigsetjmp(tm_jmpbuf, 0)
  store %setjmp_ret, tm_longjmp_ret
  store 1, tm_nested_call_counter  ; Mark outer
  call tm_begin()
  br label %cont

nested:
  %counter_inc = add %counter, 1
  store %counter_inc, tm_nested_call_counter
  br label %cont

cont:
  ; ... transaction body ...

exit:
  %counter_end = load tm_nested_call_counter
  %counter_dec = add %counter_end, -1
  store %counter_dec, tm_nested_call_counter
  %is_outer_end = icmp eq %counter_dec, 0
  br %is_outer_end, label %commit, label %return

commit:
  call tm_end()
  br label %return

return:
  ret
\end{verbatim}

This ensures:
\begin{itemize}
\item Only the outermost transaction calls \texttt{tm\_begin()}
\item Only the outermost transaction calls \texttt{tm\_end()}
\item Nested transactions inherit the outermost transaction's context
\end{enumerate}

\section{Related Work}\label{sec:related}

STM approaches fall into several categories:

\subsection{Embedded STM Languages}

\begin{itemize}
\item \textbf{Atomic Java} \cite{HerlihyMoss1993}: First STM implementation for Java
\item \textbf{STM Haskell} \cite{HarrisT2005}: STM monad in Haskell
\item \textbf{Clojure} \cite{RichHickey}: Built-in STM with ref types
\end{itemize}

\subsection{Compile-Time STM Instrumentation}

Our approach falls into this category:

\begin{itemize}
\item \textbf{TxLLVM} \cite{TxLLVM}: LLVM extension for STM (custom LLVM fork)
\item \textbf{Lock-Free STM} \cite{DiscoPanagioti}: Hardware-assisted STM
\item \textbf{RSTM} \cite{RSTM}: Multi-version STM in Java
\item \textbf{TinySTM} \cite{TinySTM}: Fast word-based STM (our primary backend)
\item \textbf{TL2} \cite{TL2}: Transactional Locking 2
\item \textbf{SwissTM} \cite{SwissTM}: Combined lazy/pessimistic STM
\item \textbf{adaptSTM} \cite{adaptSTM}: Adaptive STM tuning
\end{itemize}

\subsection{Hardware TM}

\begin{itemize}
\item \textbf{Rock} \cite{Rock}: Oracle SPARC HTM
\item \textbf{Haswell TSX} \cite{TSX}: Intel TSX
\item \textbf{IBM Power} \cite{Power}: IBM Power HTM
\end{itemize}

\subsection{STM Benchmarks}

We evaluate using established benchmarks:

\begin{itemize}
\item \textbf{STAMP} \cite{STAMP}: Stanford Transactional Memory Benchmark suite
  \begin{itemize}
  \item genome: DNA sequence matching
  \item kmeans: Clustering algorithm
  \item intruder: Intrusion detection
  \item ssca2: Graph algorithms
  \item vacation: Travel reservation system
  \item yada: Delaunay mesh
  \end{itemize}
\item \textbf{STMBench7} \cite{STMBench7}: Complex graph benchmark
  \begin{itemize}
  \item 45 operations on modules, assemblies, atomic parts
  \item Long/short traversals, structure modifications
  \end{itemize}
\item \textbf{YCSB} \cite{YCSB}: Yahoo! Cloud Serving Benchmark
\item \textbf{Bank}: Financial transactions (money transfer)
\item \textbf{Tree Benchmarks}: AVL, Red-Black trees, Hash maps
\end{itemize}

Our plugin supports all these benchmarks.

\section{Evaluation}\label{sec:evaluation}

We evaluate our plugin with the following benchmarks:

\subsection{Data Structure Benchmarks}

Results with SingleGlobalLock backend (1 thread, 10k initial, 100ms):

\begin{table}[h]
\centering
\begin{tabular}{l|r}
Benchmark & Ops/sec\\\hline
avltree & 19,081,200\\
hashmap & 3,923,730\\
bitmap & 491,876\\
list & 937,371\\
set & 3,180,490\\
\end{tabular}
\end{table}

\subsection{Bank Benchmark}

The bank benchmark transfers money between accounts:

\begin{verbatim}
Total transfers:  139,751
Total txns:       174,687
Txns/sec:        345,915
PASS: Bank total correct
\end{verbatim}

\bibliographystyle{plain}
\begin{thebibliography}{99}

%\bibitem{herlihy93}
%Maurice Herlihy and J. Eliot B. Moss. ``Transactional Memory: Architectural Support for Lock-Free Data Structures''. {\em Proceedings of ISCA}, 1993.

%\bibitem{txllvm15}
%Aleksander B. J. Christensen et al. ``TxLLVM''. \url{https://github.com/bwmania/txllvm}, 2015.

%\bibitem{tinySTM08}
%Pascal Felber et al. ``TinySTM''. \url{https://github.com/PatrickMa/STM}, 2008.

%\bibitem{tl2}
%Dave Dice, Ori Shalev, and Nir Shavit. ``Transactional Locking II''. {\em USENIX}, 2006.

%\bibitem{swissTM}
%Michal Kapalka et al. ``SwissTM: Stretching Transactional Memory''. {\em EuroSys}, 2009.

%\bibitem{adapt}
%Tiago M. R. Rico et al. ``adaptSTM: Adaptive STM''. {\em TC}, 2010.

%\bibitem{stamp}
%C. Min and S. Kim. ``STAMP Benchmark Suite''. \url{https://github.com/rmascarenes/stamp}, 2008.

%\bibitem{stampbench07}
%R. Guerraoui, M. Kapalka, and J. Vitek. ``STMBench7: A benchmark for software transactional memory''. {\em EuroSys}, 2007.

%\bibitem{ycsb}
%Yahoo! ``YCSB Benchmark''. \url{https://github.com/brianfrankcooper/YCSB}, 2010.

\end{thebibliography}

\end{document}
